% GENERATED FILE. Edit canonical lessons, then run npm run generate:books.
% canonical-source-hash: fnv1a64:c4cfe06179d230b2
% canonical-lessons: KA-C425-svikarisu, KA-C425-tiraskarisu, KA-C425-pramanikarisu, KA-C425-urjita, KA-C425-manyate, KA-R425-review-419-425

\chapter{Accepting and Rejecting}
\label{ch:ka-accepting-and-rejecting}
\hlchaptermodality{425}

\begin{chapteropening}
\textbf{By the end of this chapter:} \emph{I can say to accept, to reject, to certify, valid and recognition.}

Complete the last lesson of chapter 425: I can say to accept, to reject, to certify, valid and recognition.
\end{chapteropening}

\section[\texorpdfstring{\kn{ಸ್ವೀಕರಿಸು} (svīkarisu)}{ಸ್ವೀಕರಿಸು (svīkarisu)}]{\kn{ಸ್ವೀಕರಿಸು} (svīkarisu) — to accept; to receive}

\label{lesson:KA-C425-svikarisu}



\begin{quote}
Before the new one: say the Kannada for specific, then the Kannada for an audit.
\end{quote}

\subsection*{You'll want to know: \kn{ಸ್ವೀಕರಿಸು}}
\textbf{\kn{ಸ್ವೀಕರಿಸು}} — \emph{svīkarisu} — \textquotedblleft{}to accept; to receive\textquotedblright{}.

\begin{quote}
\kn{ಅವಳು} \kn{ಹೊಸ} \kn{ಕೆಲಸದ} \kn{ಆಹ್ವಾನವನ್ನು} \kn{ಸ್ವೀಕರಿಸಿದಳು}.
\end{quote}

\begin{quote}
\emph{avaḷu hosa kelasada āhvānavannu svīkarisidaḷu.}
\end{quote}

\begin{quote}
She accepted the offer of a new job.
\end{quote}

Past (I …): \kn{ಸ್ವೀಕರಿಸಿದೆ}, \emph{svīkariside}.

Formal. \kn{ಸ್ವೀಕಾರ}, acceptance.

\begin{cousinweb}
\textbf{\kn{ಸ್ವೀಕರಿಸು}} (\emph{svīkarisu}) is built on Sanskrit \emph{svīkāra}, \textquotedblleft{}making one's own\textquotedblright{} (Monier-Williams), with the Kannada verb ending -\emph{isu}.
\end{cousinweb}

\begin{grammarlens}[title={What you've built}]
One more word for saying what happens if.
\end{grammarlens}

\subsection*{Your turn}
\begin{itemize}
\raggedright
  \item \emph{Say it:} \emph{svīkarisu}
  \item \emph{Say it:} \emph{svīkarisu}, once more
  \item \emph{Say it:} \emph{avaḷu hosa kelasada āhvānavannu svīkarisidaḷu.} — She accepted the offer of a new job.
  \item \emph{Recall:} say the Kannada for specific, then the Kannada for an audit, then say \emph{svīkarisu} again
\end{itemize}

\begin{tcolorbox}[breakable,colback=teal!4,colframe=teal!35!black,title={Before you move on}]
What does \kn{ಸ್ವೀಕರಿಸು} mean? (\textquotedblleft{}To accept; to receive\textquotedblright{}.) Say it once more.
\end{tcolorbox}
\section[\texorpdfstring{\kn{ತಿರಸ್ಕರಿಸು} (tiraskarisu)}{ತಿರಸ್ಕರಿಸು (tiraskarisu)}]{\kn{ತಿರಸ್ಕರಿಸು} (tiraskarisu) — to reject}

\label{lesson:KA-C425-tiraskarisu}



\begin{quote}
Before the new one: say the Kannada for an audit, then the Kannada for to accept.
\end{quote}

\subsection*{You'll want to know: \kn{ತಿರಸ್ಕರಿಸು}}
\textbf{\kn{ತಿರಸ್ಕರಿಸು}} — \emph{tiraskarisu} — \textquotedblleft{}to reject\textquotedblright{}.

\begin{quote}
\kn{ಸಂಸ್ಥೆ} \kn{ನನ್ನ} \kn{ಅರ್ಜಿಯನ್ನು} \kn{ತಿರಸ್ಕರಿಸಿತು}.
\end{quote}

\begin{quote}
\emph{saṃsthe nanna arjiyannu tiraskarisitu.}
\end{quote}

\begin{quote}
The organisation rejected my application.
\end{quote}

Past (I …): \kn{ತಿರಸ್ಕರಿಸಿದೆ}, \emph{tiraskariside}.

\kn{ತಿರಸ್ಕಾರ}, contempt or rejection.

\begin{cousinweb}
\textbf{\kn{ತಿರಸ್ಕರಿಸು}} (\emph{tiraskarisu}) is built on Sanskrit \emph{tiraskāra}, \textquotedblleft{}placing aside; contempt\textquotedblright{} (Monier-Williams), with the Kannada verb ending -\emph{isu}.
\end{cousinweb}

\begin{grammarlens}[title={What you've built}]
One more word for saying what happens if.
\end{grammarlens}

\subsection*{Your turn}
\begin{itemize}
\raggedright
  \item \emph{Say it:} \emph{tiraskarisu}
  \item \emph{Say it:} \emph{tiraskarisu}, once more
  \item \emph{Say it:} \emph{saṃsthe nanna arjiyannu tiraskarisitu.} — The organisation rejected my application.
  \item \emph{Recall:} say the Kannada for an audit, then the Kannada for to accept, then say \emph{tiraskarisu} again
\end{itemize}

\begin{tcolorbox}[breakable,colback=teal!4,colframe=teal!35!black,title={Before you move on}]
What does \kn{ತಿರಸ್ಕರಿಸು} mean? (\textquotedblleft{}To reject\textquotedblright{}.) Say it once more.
\end{tcolorbox}
\section[\texorpdfstring{\kn{ಪ್ರಮಾಣೀಕರಿಸು} (pramāṇīkarisu)}{ಪ್ರಮಾಣೀಕರಿಸು (pramāṇīkarisu)}]{\kn{ಪ್ರಮಾಣೀಕರಿಸು} (pramāṇīkarisu) — to certify}

\label{lesson:KA-C425-pramanikarisu}



\begin{quote}
Before the new one: say the Kannada for to accept, then the Kannada for to reject.
\end{quote}

\subsection*{You'll want to know: \kn{ಪ್ರಮಾಣೀಕರಿಸು}}
\textbf{\kn{ಪ್ರಮಾಣೀಕರಿಸು}} — \emph{pramāṇīkarisu} — \textquotedblleft{}to certify\textquotedblright{}.

\begin{quote}
\kn{ವೈದ್ಯರು} \kn{ನನ್ನ} \kn{ಆರೋಗ್ಯವನ್ನು} \kn{ಪ್ರಮಾಣೀಕರಿಸಿದರು}.
\end{quote}

\begin{quote}
\emph{vaidyaru nanna ārōgyavannu pramāṇīkarisidaru.}
\end{quote}

\begin{quote}
The doctor certified my health.
\end{quote}

Past (I …): \kn{ಪ್ರಮಾಣೀಕರಿಸಿದೆ}, \emph{pramāṇīkariside}.

\kn{ಪ್ರಮಾಣಪತ್ರ}, a certificate, is the paper you get.

\begin{cousinweb}
\textbf{\kn{ಪ್ರಮಾಣೀಕರಿಸು}} (\emph{pramāṇīkarisu}) is built on Sanskrit \emph{pramāṇa}, \textquotedblleft{}a standard, a proof\textquotedblright{} (CDIAL 8723), with -\emph{īkarisu}, \textquotedblleft{}to make into\textquotedblright{}: to make proof of.
\end{cousinweb}

\begin{grammarlens}[title={What you've built}]
One more word for saying what happens if.
\end{grammarlens}

\subsection*{Your turn}
\begin{itemize}
\raggedright
  \item \emph{Say it:} \emph{pramāṇīkarisu}
  \item \emph{Say it:} \emph{pramāṇīkarisu}, once more
  \item \emph{Say it:} \emph{vaidyaru nanna ārōgyavannu pramāṇīkarisidaru.} — The doctor certified my health.
  \item \emph{Recall:} say the Kannada for to accept, then the Kannada for to reject, then say \emph{pramāṇīkarisu} again
\end{itemize}

\begin{tcolorbox}[breakable,colback=teal!4,colframe=teal!35!black,title={Before you move on}]
What does \kn{ಪ್ರಮಾಣೀಕರಿಸು} mean? (\textquotedblleft{}To certify\textquotedblright{}.) Say it once more.
\end{tcolorbox}
\section[\texorpdfstring{\kn{ಊರ್ಜಿತ} (ūrjita)}{ಊರ್ಜಿತ (ūrjita)}]{\kn{ಊರ್ಜಿತ} (ūrjita) — valid; in force}

\label{lesson:KA-C425-urjita}



\begin{quote}
Before the new one: say the Kannada for to reject, then the Kannada for to certify.
\end{quote}

\subsection*{You'll want to know: \kn{ಊರ್ಜಿತ}}
\textbf{\kn{ಊರ್ಜಿತ}} — \emph{ūrjita} — \textquotedblleft{}valid; in force\textquotedblright{}.

\begin{quote}
\kn{ಈ} \kn{ಪಾಸ್} \kn{ಇನ್ನೂ} \kn{ಒಂದು} \kn{ತಿಂಗಳು} \kn{ಊರ್ಜಿತ}.
\end{quote}

\begin{quote}
\emph{ī pās innū ondu tiṅgaḷu ūrjita.}
\end{quote}

\begin{quote}
This pass is valid for one more month.
\end{quote}

\kn{ಊರ್ಜಿತವಾಗು}, to stay in force.

\begin{cousinweb}
\textbf{\kn{ಊರ್ಜಿತ}} (\emph{ūrjita}) is borrowed from Sanskrit \emph{ūrjita}, \textquotedblleft{}endowed with strength, powerful\textquotedblright{} (Monier-Williams).
\end{cousinweb}

\begin{grammarlens}[title={What you've built}]
One more word for saying what happens if.
\end{grammarlens}

\subsection*{Your turn}
\begin{itemize}
\raggedright
  \item \emph{Say it:} \emph{ūrjita}
  \item \emph{Say it:} \emph{ūrjita}, once more
  \item \emph{Say it:} \emph{ī pās innū ondu tiṅgaḷu ūrjita.} — This pass is valid for one more month.
  \item \emph{Recall:} say the Kannada for to reject, then the Kannada for to certify, then say \emph{ūrjita} again
\end{itemize}

\begin{tcolorbox}[breakable,colback=teal!4,colframe=teal!35!black,title={Before you move on}]
What does \kn{ಊರ್ಜಿತ} mean? (\textquotedblleft{}Valid; in force\textquotedblright{}.) Say it once more.
\end{tcolorbox}
\section[\texorpdfstring{\kn{ಮಾನ್ಯತೆ} (mānyate)}{ಮಾನ್ಯತೆ (mānyate)}]{\kn{ಮಾನ್ಯತೆ} (mānyate) — recognition; validity}

\label{lesson:KA-C425-manyate}



\begin{quote}
Before the new one: say the Kannada for to certify, then the Kannada for valid.
\end{quote}

\subsection*{You'll want to know: \kn{ಮಾನ್ಯತೆ}}
\textbf{\kn{ಮಾನ್ಯತೆ}} — \emph{mānyate} — \textquotedblleft{}recognition; validity\textquotedblright{}.

\begin{quote}
\kn{ಈ} \kn{ಕಾಲೇಜಿಗೆ} \kn{ಸರ್ಕಾರದ} \kn{ಮಾನ್ಯತೆ} \kn{ಇದೆ}.
\end{quote}

\begin{quote}
\emph{ī kālējige sarkārada mānyate ide.}
\end{quote}

\begin{quote}
This college is recognised by the government.
\end{quote}

\kn{ಮಾನ್ಯ}, recognised, honoured.

\begin{cousinweb}
\textbf{\kn{ಮಾನ್ಯತೆ}} (\emph{mānyate}) is built on Sanskrit \emph{mānya}, \textquotedblleft{}to be respected or honoured\textquotedblright{} (Monier-Williams), with -\emph{tā}, \textquotedblleft{}-ness\textquotedblright{}.
\end{cousinweb}

\begin{grammarlens}[title={What you've built}]
One more word for saying what happens if.
\end{grammarlens}

\subsection*{Your turn}
\begin{itemize}
\raggedright
  \item \emph{Say it:} \emph{mānyate}
  \item \emph{Say it:} \emph{mānyate}, once more
  \item \emph{Say it:} \emph{ī kālējige sarkārada mānyate ide.} — This college is recognised by the government.
  \item \emph{Recall:} say the Kannada for to certify, then the Kannada for valid, then say \emph{mānyate} again
\end{itemize}

\begin{tcolorbox}[breakable,colback=teal!4,colframe=teal!35!black,title={Before you move on}]
What does \kn{ಮಾನ್ಯತೆ} mean? (\textquotedblleft{}Recognition; validity\textquotedblright{}.) Say it once more.
\end{tcolorbox}
\section[(dialogue)]{Looking back: from Success and Failure to Accepting and Rejecting}

\label{lesson:KA-R425-review-419-425}



\begin{quote}
Say the Kannada for to fall through and to bear fruit.
\end{quote}

\subsection*{Your turn}
Say these aloud:

\begin{itemize}
\raggedright
  \item to fall through, to bear fruit, to come true, a crisis, unavoidable
  \item to waver, to wake up to, compulsion, uncertain, a calamity
  \item a demand, to supply, a shortage, export, inflation
  \item to refer to, a criterion, a course of action, a method, a process
  \item a stage, a task, a strategy, a trick, to divide
  \item to observe, to critique, a circular, specific, an audit
  \item to accept, to reject, to certify, valid, recognition
\end{itemize}

\begin{tcolorbox}[breakable,colback=teal!4,colframe=teal!35!black,title={Before you move on}]
To fall through? (\textbf{\kn{ವಿಫಲವಾಗು}}, \emph{viphalavāgu}.) To bear fruit? (\textbf{\kn{ಫಲಿಸು}}, \emph{phalisu}.) To come true? (\textbf{\kn{ಕೈಗೂಡು}}, \emph{kaigūḍu}.) A crisis? (\textbf{\kn{ಬಿಕ್ಕಟ್ಟು}}, \emph{bikkaṭṭu}.) Unavoidable? (\textbf{\kn{ಅನಿವಾರ್ಯ}}, \emph{anivārya}.) To waver? (\textbf{\kn{ಅಳುಕು}}, \emph{aḷuku}.) To wake up to? (\textbf{\kn{ಎಚ್ಚೆತ್ತುಕೊ}}, \emph{eccettuko}.) Compulsion? (\textbf{\kn{ಬಲವಂತ}}, \emph{balavanta}.) Uncertain? (\textbf{\kn{ಅನಿಶ್ಚಿತ}}, \emph{aniścita}.) A calamity? (\textbf{\kn{ಆಪತ್ತು}}, \emph{āpattu}.) A demand? (\textbf{\kn{ಬೇಡಿಕೆ}}, \emph{bēḍike}.) To supply? (\textbf{\kn{ಪೂರೈಸು}}, \emph{pūraisu}.) A shortage? (\textbf{\kn{ಕೊರತೆ}}, \emph{korate}.) Export? (\textbf{\kn{ರಫ್ತು}}, \emph{raphtu}.) Inflation? (\textbf{\kn{ಹಣದುಬ್ಬರ}}, \emph{haṇadubbara}.) To refer to? (\textbf{\kn{ಪರಾಮರ್ಶಿಸು}}, \emph{parāmarśisu}.) A criterion? (\textbf{\kn{ಮಾನದಂಡ}}, \emph{mānadaṇḍa}.) A course of action? (\textbf{\kn{ಕ್ರಮ}}, \emph{krama}.) A method? (\textbf{\kn{ವಿಧಾನ}}, \emph{vidhāna}.) A process? (\textbf{\kn{ಪ್ರಕ್ರಿಯೆ}}, \emph{prakriye}.) A stage? (\textbf{\kn{ಹಂತ}}, \emph{hanta}.) A task? (\textbf{\kn{ಕಾರ್ಯ}}, \emph{kārya}.) A strategy? (\textbf{\kn{ತಂತ್ರ}}, \emph{tantra}.) A trick? (\textbf{\kn{ಉಪಾಯ}}, \emph{upāya}.) To divide? (\textbf{\kn{ವಿಭಜಿಸು}}, \emph{vibhajisu}.) To observe? (\textbf{\kn{ಅವಲೋಕಿಸು}}, \emph{avalōkisu}.) To critique? (\textbf{\kn{ವಿಮರ್ಶಿಸು}}, \emph{vimarśisu}.) A circular? (\textbf{\kn{ಸುತ್ತೋಲೆ}}, \emph{suttōle}.) Specific? (\textbf{\kn{ನಿರ್ದಿಷ್ಟ}}, \emph{nirdiṣṭa}.) An audit? (\textbf{\kn{ಪರಿಶೋಧನೆ}}, \emph{pariśōdhane}.) To accept? (\textbf{\kn{ಸ್ವೀಕರಿಸು}}, \emph{svīkarisu}.) To reject? (\textbf{\kn{ತಿರಸ್ಕರಿಸು}}, \emph{tiraskarisu}.) To certify? (\textbf{\kn{ಪ್ರಮಾಣೀಕರಿಸು}}, \emph{pramāṇīkarisu}.) Valid? (\textbf{\kn{ಊರ್ಜಿತ}}, \emph{ūrjita}.) Recognition? (\textbf{\kn{ಮಾನ್ಯತೆ}}, \emph{mānyate}.)
\end{tcolorbox}
